%% SHARED SECTION -- the proof of Proposition~\ref{prop:lmmse}, included from
%% the appendix. Split out of prop-lmmse.tex because the body has a nine-page
%% limit and a full induction does not belong in it; the STATEMENT, the
%% projection definition and the hypothesis discussion stay in the body, which
%% is what a reader needs there.

\begin{remark}[On the hypotheses]
\label{rem:lmmse-hyp}
Four are load-bearing and worth separating. \emph{Finite second moments} are what make $\Pgauss$
defined at all. \emph{Zero mean} is not cosmetic: the general LMMSE carries
$\E[\va] + \Cov(\va,\vx)\Cov(\vx)^{-1}(\vx-\E[\vx])$, and \eqref{eq:lmmse} is the centred
special case. \emph{$\corr\neq0$} is required because the backward map in \eqref{eq:closure}
divides by $\corr$; at $\corr=0$ the chain has no edges, the sites are independent, and the
claim holds trivially and separately, each site's posterior mean being the scalar Gaussian
$\Sig_0(\Sig_0+\Dt e^{2t})^{-1}x_i$. \emph{Where the projection is applied} matters because
$\Pgauss$ does not commute with multiplication by the likelihood; Proposition~\ref{prop:lmmse}
fixes the order, and a different order is a different algorithm.

The statement is about exact message functions on $\Reals$ and an analytic projection operator.
\textbf{The grid plays no part in it.} Discretisation enters only when the \emph{unrestricted}
non-Gaussian messages are represented numerically, which is what the recursion is compared
against; conflating the two would turn a theorem about a closure into a claim about quadrature.
\end{remark}

\begin{proof}
Write $\mu_i^{\rightarrow}$ for the forward message at site $i$ after projection and
$\ell_i$ for the Gaussian unary likelihood. We show by induction that every projected message
coincides with the corresponding exact message of the covariance-matched Gaussian chain --- the
chain with the same $\corr$ and innovation variance $\ivar$ but Gaussian innovations.

\emph{Base.} The initial message is Gaussian by hypothesis, or is projected to the Gaussian with
the same first two moments; either way it equals the Gaussian chain's initial message, which has
those same moments by construction.

\emph{Step.} Suppose $\mu_{i-1}^{\rightarrow}$ agrees with the Gaussian chain's message, so it is
Gaussian with some mean $m$ and variance $v$. Multiplying by $\ell_{i-1}$ multiplies two Gaussian
densities and yields a Gaussian, with mean and variance given by the usual precision-weighted
combination --- identical in both chains, since the likelihood is the same object. Taking the
transition integral against $a_i = \corr a_{i-1} + \varepsilon_i$ convolves with $\innov$ and
scales by $\corr$, which maps $(m,v) \mapsto (\corr m,\, \corr^2 v + \ivar)$: the mean and variance
of a convolution depend on $\innov$ \emph{only} through its mean and variance, so this is the same
map for every innovation with moments $(0,\ivar)$, and in particular the same as for the Gaussian
chain. In the Gaussian chain the convolution of two Gaussians is already Gaussian, so no
projection occurs there; here $\Pgauss$ is applied, and since the pre-projection function has
exactly the moments just computed, the projection returns precisely the Gaussian chain's message.
The backward recursion is identical with the transition map inverted, $(m,v)\mapsto
(m/\corr,\,(v+\ivar)/\corr^2)$, which is where $\corr\neq0$ is used.

Hence every projected belief equals the corresponding Gaussian-chain belief, and the returned
posterior mean is the Gaussian chain's posterior mean. For jointly Gaussian zero-mean
$(\va,\vx)$ the conditional mean is linear in $\vx$ and therefore coincides with the best linear
estimator, which is \eqref{eq:lmmse}: since $\vx = e^{-t}\va + \sqrt{\Dt}\,\vect{z}$ with
$\vect{z}$ independent, $\Cov(\va,\vx)=e^{-t}\Sig_0$ and $\Cov(\vx)=e^{-2t}\Sig_0+\Dt\Id$. Since
the recursion never read $\innov$ beyond $(0,\ivar)$, the same estimator is returned for every
innovation law with those moments.
\end{proof}
